% ACI Companion — Lesson 4: Evolutionary Algorithms — GA & ACO
\documentclass[a4paper,11pt]{article}
\usepackage[T1]{fontenc}
\usepackage[utf8]{inputenc}
\usepackage{lmodern}
\usepackage{microtype}
\usepackage[a4paper,margin=1in,headheight=15pt,headsep=14pt]{geometry}
\usepackage{amsmath,amssymb}
\usepackage{graphicx}
\usepackage{booktabs}
\usepackage[table,svgnames]{xcolor}
\usepackage[most]{tcolorbox}
\tcbuselibrary{skins,breakable}
\usepackage{pifont}
\usepackage{enumitem}
\setlist{leftmargin=1.5em,topsep=4pt,itemsep=3pt}
\usepackage{parskip}
\usepackage{fancyhdr}
\usepackage{titlesec}
\usepackage[hidelinks]{hyperref}
\usepackage{xurl}

\definecolor{NavyBanner}{HTML}{21355E}
\definecolor{IntFrame}{HTML}{2C5AA0}
\definecolor{IntTint}{HTML}{EBF0FA}
\definecolor{EveFrame}{HTML}{8A5A1E}
\definecolor{EveTint}{HTML}{FDF3E7}
\definecolor{WrkFrame}{HTML}{2E7D52}
\definecolor{WrkTint}{HTML}{EBF5EF}
\definecolor{WatFrame}{HTML}{B23A48}
\definecolor{WatTint}{HTML}{FAE9EB}
\definecolor{KeyFrame}{HTML}{6A4C93}
\definecolor{KeyTint}{HTML}{F0EBF8}
\definecolor{SectionNavy}{HTML}{1A2744}

\newtcolorbox{intuition}{enhanced,breakable,colback=IntTint,colframe=IntFrame,
  boxrule=0.6pt,arc=4pt,attach boxed title to top left={xshift=6mm,yshift=-3mm},
  boxed title style={colback=IntFrame,rounded corners},coltitle=white,
  fonttitle=\bfseries\sffamily\small,title={\ding{43}\; The intuition}}
\newtcolorbox{everyday}{enhanced,breakable,colback=EveTint,colframe=EveFrame,
  boxrule=0.6pt,arc=4pt,attach boxed title to top left={xshift=6mm,yshift=-3mm},
  boxed title style={colback=EveFrame,rounded corners},coltitle=white,
  fonttitle=\bfseries\sffamily\small,title={\ding{72}\; Everyday picture}}
\newtcolorbox[auto counter]{worked}[1]{enhanced,breakable,colback=WrkTint,colframe=WrkFrame,
  boxrule=0.6pt,arc=4pt,attach boxed title to top left={xshift=6mm,yshift=-3mm},
  boxed title style={colback=WrkFrame,rounded corners},coltitle=white,
  fonttitle=\bfseries\sffamily\small,title={\ding{46}\; Worked example: #1}}
\newtcolorbox{watchout}{enhanced,breakable,colback=WatTint,colframe=WatFrame,
  boxrule=0.6pt,arc=4pt,attach boxed title to top left={xshift=6mm,yshift=-3mm},
  boxed title style={colback=WatFrame,rounded corners},coltitle=white,
  fonttitle=\bfseries\sffamily\small,title={\ding{55}\; Watch out}}
\newtcolorbox{keytake}{enhanced,breakable,colback=KeyTint,colframe=KeyFrame,
  boxrule=0.6pt,arc=4pt,attach boxed title to top left={xshift=6mm,yshift=-3mm},
  boxed title style={colback=KeyFrame,rounded corners},coltitle=white,
  fonttitle=\bfseries\sffamily\small,title={\ding{52}\; Key takeaway}}

\titleformat{\section}{\large\bfseries\sffamily\color{SectionNavy}}{}{0em}{}[\titlerule]
\titleformat{\subsection}{\normalsize\bfseries\sffamily\color{SectionNavy}}{}{0em}{}
\pagestyle{fancy}\fancyhf{}
\lhead{\small\textsf{ACI Companion}}
\rhead{\small\textsf{Lesson 4 · Evolutionary Algorithms: GA \& ACO}}
\renewcommand{\headrulewidth}{0.4pt}
\cfoot{\small\thepage}

\begin{document}

\begin{tcolorbox}[enhanced,colback=NavyBanner,colframe=NavyBanner,arc=0pt,
  boxrule=0pt,left=10mm,right=10mm,top=8mm,bottom=8mm]
  {\small\bfseries\sffamily\color{white}\MakeUppercase{ACI · AIMLCZG557}}\par\vspace{4pt}
  {\LARGE\bfseries\sffamily\color{white}Evolutionary Algorithms: GA \& ACO}\par\vspace{4pt}
  {\normalsize\color{white!80!NavyBanner}Advanced Computing for AI — Lesson 4}\par
  \textcolor{white!40!NavyBanner}{\rule{\linewidth}{0.4pt}}\par\vspace{2pt}
  {\small\color{white!70!NavyBanner}Exam: 28 Jun 2026 · ACO pheromone formula + GA generation trace are key exam items}
\end{tcolorbox}

\vspace{4pt}
{\small\textbf{How to use this companion.} \textcolor{IntFrame}{Blue} = the idea. \textcolor{EveFrame}{Amber} = everyday analogy. \textcolor{WrkFrame}{Green} = fully-solved example. \textcolor{WatFrame}{Red} = common mistake. \textcolor{KeyFrame}{Purple} = one line to remember.}

\vspace{8pt}

%=============================================================================
\section{1 \quad Genetic Algorithms — The Big Picture}

A Genetic Algorithm (GA) mimics natural evolution. You maintain a \textbf{population} of candidate solutions (called \textbf{chromosomes}). Each chromosome is scored by a \textbf{fitness function}. Better-scoring individuals are more likely to be selected as parents. Parents combine via \textbf{crossover} to produce offspring. Random \textbf{mutation} maintains diversity. Over many generations, the population evolves toward better solutions.

\begin{everyday}
Think of dog breeding. You pick the fastest dogs (selection), mate them (crossover), occasionally a puppy gets an unexpected trait (mutation), and after many generations you have greyhounds. A GA does the same, but the ``dogs'' are encoded solutions to a problem.
\end{everyday}

\textbf{GA Pseudocode (exam-ready):}
\begin{verbatim}
function GA(population, fitness):
  evaluate all individuals
  while not termination_condition:
    parents = SELECT(population, fitness)
    children = CROSSOVER(parents)
    children = MUTATE(children)
    evaluate children
    population = REPLACE(population, children)
  return best_individual_found
\end{verbatim}

\textbf{Where this shows up in ML.} GAs are used in Neural Architecture Search (NAS), hyperparameter optimisation, and evolutionary strategies for RL. Google's AutoML-Zero used GAs to discover entire ML algorithms from scratch.

\begin{keytake}
GA = population + selection + crossover + mutation + replacement. Each iteration is one generation. The fitness function is the only problem-specific part.
\end{keytake}

%=============================================================================
\section{2 \quad Selection Methods}

\textbf{Tournament selection.} Pick $k$ individuals at random from the population. The one with the highest fitness wins and becomes a parent. Common: $k=2$ (binary tournament). Higher $k$ = more selection pressure (less diversity).

\textbf{Roulette wheel (fitness-proportionate).} Each individual gets a slice of a wheel proportional to its fitness. Selection probability of individual $i$:
\[
P(i) = \frac{f_i}{\sum_{j=1}^{N} f_j}
\]
Problem: if one individual dominates, it gets nearly all the wheel (premature convergence).

\textbf{Rank-based.} Sort by fitness; assign probability based on rank, not raw fitness. Prevents one dominant individual from taking over.

\begin{worked}{Roulette wheel probabilities}
Population of 4: $f_1=10,\ f_2=5,\ f_3=3,\ f_4=2$. Total = 20.
\begin{enumerate}[label=\textbf{Step \arabic*.}]
  \item $P(1) = 10/20 = 0.50$
  \item $P(2) = 5/20 = 0.25$
  \item $P(3) = 3/20 = 0.15$
  \item $P(4) = 2/20 = 0.10$
\end{enumerate}
Individual 1 wins half the time. If it is only slightly better than the others, this may not be ideal — use rank-based selection instead.
\end{worked}

\begin{keytake}
Tournament = simple, no sorting, controllable pressure. Roulette = proportional but risky. Rank-based = robust to fitness scaling.
\end{keytake}

%=============================================================================
\section{3 \quad Crossover Operators}

Crossover combines two parent chromosomes to create offspring.

\textbf{Single-point crossover.} Pick a random cut point. Child 1 = left part of Parent A + right part of Parent B. Child 2 = left part of Parent B + right part of Parent A.

\textbf{Two-point crossover.} Pick two cut points. Swap the middle segment.

\textbf{Uniform crossover.} For each gene position, flip a coin. Child inherits from Parent A with probability 0.5, else Parent B.

\begin{worked}{Single-point crossover, step by step}
Parent A: \texttt{1 1 0 0 1 0 1} \quad Parent B: \texttt{0 0 1 1 0 1 0} \quad Cut point: position 3.
\begin{enumerate}[label=\textbf{Step \arabic*.}]
  \item Child 1 = first 3 genes from A + last 4 from B: \texttt{1 1 0 | 1 0 1 0}
  \item Child 2 = first 3 genes from B + last 4 from A: \texttt{0 0 1 | 0 1 0 1}
\end{enumerate}
Both children inherit traits from both parents. The crossover point is chosen uniformly at random each generation.
\end{worked}

%=============================================================================
\section{4 \quad Mutation}

Mutation applies a small random change to an offspring's chromosome. For binary chromosomes, flip each bit with probability $p_m$ (typically $1/L$ where $L$ is chromosome length).

\begin{intuition}
Mutation prevents the population from getting stuck. Without mutation, once all individuals share the same allele at a position, crossover can never reintroduce diversity there. Mutation is the only source of truly new genetic material.
\end{intuition}

\begin{watchout}
Mutation rate $p_m$ is a key hyperparameter. Too high: random walk (no learning). Too low: premature convergence. Typical range: $0.001$ to $0.01$ per bit.
\end{watchout}

%=============================================================================
\section{5 \quad Generational vs Steady-State GA}

This is a \textbf{confirmed exam error} — make sure you know the difference.

\textbf{Generational GA.} The entire population is replaced at each generation. All $N$ parents produce $N$ children; the new population consists entirely of children. Generation boundary is clear.

\textbf{Steady-state GA.} Only a \emph{few} individuals are replaced at a time (often just 1 or 2). The rest of the population survives unchanged. The population evolves incrementally. This maintains more diversity and is more suitable for online/continuous scenarios.

\begin{watchout}
The most common exam mistake: saying steady-state GA ``replaces the entire population every generation''. That is WRONG. Steady-state replaces only a few individuals. Only generational GA replaces the whole population. This distinction appeared in your lesson as a \textbf{confirmed error to fix}.
\end{watchout}

\begin{keytake}
Generational = replace ALL $N$ individuals each generation. Steady-state = replace only 1--2 individuals at a time. The exam will test this.
\end{keytake}

%=============================================================================
\section{6 \quad Crowding \& Diversity Maintenance}

When many individuals converge on the same region, the population loses diversity (premature convergence). Crowding methods maintain diversity by preventing similar individuals from dominating.

\textbf{Fitness sharing.} Reduce the fitness of an individual based on how many similar individuals are nearby. If many individuals cluster around a peak, each gets a smaller share of that peak's fitness. Formula:
\[
f'_i = \frac{f_i}{\sum_{j=1}^{N} \mathrm{sh}(d_{ij})}
\]
where $\mathrm{sh}(d) = 1 - (d/\sigma_{\text{share}})^\alpha$ if $d < \sigma_{\text{share}}$, else 0. $d_{ij}$ is the distance between individuals $i$ and $j$.

\textbf{Deterministic crowding.} Each child replaces the most similar parent, keeping the population diverse without explicit distance calculations.

\begin{watchout}
A diversity-maximising replacement strategy is NOT the same as generational replacement. Crowding methods (fitness sharing, crowding distance, deterministic crowding) specifically exist to \emph{counteract} the tendency of selection to collapse diversity.
\end{watchout}

\textbf{Where this shows up in ML.} Quality-Diversity (QD) algorithms in RL (like MAP-Elites) use fitness sharing ideas to maintain a diverse archive of behaviours, not just a single best policy.

%=============================================================================
\section{7 \quad Ant Colony Optimisation (ACO)}

ACO mimics how real ants find the shortest path to food. Ants deposit \textbf{pheromone} on paths they travel. Shorter paths accumulate more pheromone per unit time. Other ants preferentially follow high-pheromone paths. Over time, the shortest path gets reinforced.

\begin{intuition}
Imagine a colony of ants exploring between nest and food. A short path is completed faster, so ants on it deposit pheromone more frequently. More pheromone attracts more ants, which deposit more pheromone. It is a positive feedback loop that amplifies the best solution.
\end{intuition}

\subsection{7.1 Edge Probability}

An ant at node $i$ chooses the next node $j$ with probability:
\[
p_{ij} = \frac{\tau_{ij}^\alpha \cdot \eta_{ij}^\beta}{\displaystyle\sum_{k \in \text{allowed}} \tau_{ik}^\alpha \cdot \eta_{ik}^\beta}
\]

\begin{itemize}
  \item $\tau_{ij}$ — pheromone level on edge $(i,j)$. Higher = more attractive (learned).
  \item $\eta_{ij} = 1/d_{ij}$ — heuristic (inverse distance). Higher = shorter edge (domain knowledge).
  \item $\alpha$ — pheromone importance. $\alpha = 0$: ignore pheromone (pure greedy heuristic).
  \item $\beta$ — heuristic importance. $\beta = 0$: ignore heuristic (pure pheromone following).
\end{itemize}

\subsection{7.2 Pheromone Update}

After all ants complete their tours, pheromones are updated:
\[
\tau_{ij} \leftarrow (1 - \rho)\,\tau_{ij} + \Delta\tau_{ij}
\]

\begin{itemize}
  \item $(1-\rho)\,\tau_{ij}$ — \textbf{evaporation}: reduces all pheromones by factor $(1-\rho)$ to allow forgetting old paths. $\rho \in (0,1)$ is the evaporation rate.
  \item $\Delta\tau_{ij} = \sum_{k=1}^{m} \Delta\tau_{ij}^k$ — total new pheromone deposited by all ants on edge $(i,j)$ this iteration.
  \item $\Delta\tau_{ij}^k = Q/L_k$ if ant $k$ used edge $(i,j)$; else 0. Here $Q$ is a constant and $L_k$ is the tour length.
\end{itemize}

\begin{worked}{ACO pheromone update, one iteration}
Two ants, one edge $(A,B)$. $\tau_{AB} = 0.8$, $\rho = 0.3$, $Q = 1$.
Ant 1 tour length $L_1 = 10$ (used edge $AB$). Ant 2 tour length $L_2 = 5$ (did NOT use $AB$).

\begin{enumerate}[label=\textbf{Step \arabic*.}]
  \item Evaporate: $(1 - 0.3) \times 0.8 = 0.7 \times 0.8 = 0.56$.
  \item Deposit from ant 1: $\Delta\tau^1_{AB} = 1/10 = 0.10$.
  \item Deposit from ant 2: $\Delta\tau^2_{AB} = 0$ (ant 2 didn't use $AB$).
  \item New pheromone: $\tau_{AB} = 0.56 + 0.10 = \mathbf{0.66}$.
\end{enumerate}
The edge pheromone dropped because ant 1 had a long tour. If ant 1 had $L_1 = 2$, deposit would be $0.50$ and $\tau_{AB}$ would be $0.56 + 0.50 = 1.06$ — rewarding the shorter tour.
\end{worked}

\begin{watchout}
Evaporation is applied \emph{before} adding new pheromone. Without evaporation, pheromone accumulates forever and the system converges prematurely to whatever was good early on. Evaporation is the forgetting mechanism.
\end{watchout}

\textbf{Where this shows up in ML.} ACO has been applied to feature selection, neural architecture search, and combinatorial optimisation (routing, scheduling). Its key insight — positive feedback reinforcing good solutions — is related to bandit algorithms and Thompson sampling.

\begin{keytake}
ACO: pheromone $\times$ heuristic drives edge probability. Update: evaporate first, then deposit inversely proportional to tour length. Short tours = more pheromone = future ants follow.
\end{keytake}

\end{document}
